\chapter{Classes and Objects}

\section{Introduction}

\begin{tcolorbox}[title=Goal]
\textbf{Refactor} your Session 5 functional pipeline into a class (\texttt{IrisRuleClassifier}) that bundles configuration (threshold and labels) with behavior (prediction and evaluation). 

\textbf{What is Object-Oriented Programming (OOP)?}
\begin{itemize}\setlength{\itemsep}{2pt}
    \item A \textbf{class} is a blueprint that describes what data and behavior a thing should have.
    \item An \textbf{object} is a single instance created from that blueprint.
    \item An \textbf{attribute} is data stored on the object (for example, \texttt{self.threshold}).
    \item A \textbf{method} is a function stored on the class that operates on the object (for example, \texttt{predict()} and \texttt{evaluate()}).
\end{itemize}
\end{tcolorbox}

\begin{tcolorbox}[title=You will work in]
Use the following file for this session:
\begin{itemize}
    \item \texttt{session6/session6\_iris\_template.py} (main student file)
\end{itemize}
\textbf{How to run the file:} Always run this script from the project root using the terminal command: \\
\texttt{python3 session6/session6\_iris\_template.py} \\
\textbf{Do not edit files outside} \texttt{session6/}.
\end{tcolorbox}

In session 4 and session 5, you built an Iris classifier using \textbf{functions}.
Your pipeline worked, but you may have noticed a design pattern:
many functions needed the same configuration values (for example, the threshold and the label names), so you passed a dictionary called \texttt{settings} around repeatedly.

\noindent
\textbf{Why refactor?}
That function-based design is fine for one small classifier, but it becomes awkward as soon as you want to run experiments:
\begin{itemize}\setlength{\itemsep}{2pt}
    \item If you want 10 classifiers with 10 thresholds, you now need 10 different \texttt{settings} dictionaries.
    \item Every time you call a function, you must remember to pass the correct \texttt{settings}.
    \item If you add one new setting later, you must update many function signatures and many function calls.
\end{itemize}

\noindent
This session (Session 6), you refactor your session 5 design using \textbf{Object-Oriented Programming (OOP)}.
The core idea is simple:
\begin{quote}
\textbf{Bundle the data (configuration) and the behavior (functions) together.}
\end{quote}

\noindent
\textbf{Refactoring promise:} the classifier rule itself stays the same.
This session reorganizes code, but it should \textbf{not change} what "setosa" means or how the threshold rule works.

\subsection{Transition Map: Session 5 Functions $\rightarrow$ Session 6 Class}
The classifier logic itself (if petal\_length < threshold ...) remains exactly the same.
What changes is \textit{where} that logic lives and \textit{how} you access it.

\noindent
Use this table as your translation guide. It shows how session 5 functions and settings translate into session 6 attributes and methods:

\begin{landscape}
\begin{table}[H]
\centering
\small
\renewcommand{\arraystretch}{0.3}

\begin{tabular}{|p{0.30\textwidth}|p{0.30\textwidth}|p{0.35\textwidth}|}
\hline
\textbf{Session 5 (Functions + Settings Dict)} & \textbf{Session 6 (Class + Self State)} & \textbf{Explanation of Change} \\
\hline
\texttt{get\_default\_settings()} returns a dict: \texttt{\{"threshold": 2.0, \dots\}} & \texttt{\_\_init\_\_(self, threshold=2.0)} & Instead of passing a dictionary around, we store \texttt{self.threshold} inside the object when it is created. \\
\hline
\texttt{determine\_binary\_label(sample, settings)} & \texttt{predict(self, sample)} & The function now belongs to the class. It uses \texttt{self.threshold} instead of \texttt{settings["threshold"]}. \\
\hline
\texttt{determine\_true\_binary\_label(sample, settings)} & \textit{(Helper method or part of evaluate)} & Typically logic that calculates "truth" can be a method like \texttt{\_get\_true\_label(self, sample)} or inside evaluation. \\
\hline
\texttt{initialize\_predictions(dataset, settings, print\_each)} & \texttt{evaluate(self, dataset, print\_each)} & The evaluation loop is now a method of the classifier. It iterates over data using its own prediction logic. \\
\hline
\texttt{setup\_application(filepath, ...)} & \texttt{main()} block creating an instance & We replace the procedural setup function with object instantiation: \texttt{clf = IrisRuleClassifier(...)}. \\
\hline
\end{tabular}
\caption{Refactoring Map: From Session 5 Functional Style to Session 6 Object-Oriented Style.}
\end{table}
\end{landscape}


\section{The Blueprint and Data (\texttt{\_\_init\_\_})}

\begin{tcolorbox}[title=Do: Define the Blueprint and Data (\texttt{\_\_init\_\_})]
A class starts with a special constructor method called \texttt{\_\_init\_\_}.
This method runs \textbf{automatically} when you create a new object.
Its job is to put the object into a valid starting state.

In session 5, you used \texttt{get\_default\_settings()} to build a dictionary.
In session 6, you store that information as \textbf{attributes} on \texttt{self}.
\texttt{self} is how a method refers to the current object.

\begin{minted}[fontsize=\small, linenos, frame=lines]{python}
class IrisRuleClassifier:
    def __init__(self, threshold=2.0):
        # Store configuration inside the instance (self)
        self.threshold = threshold
        
        # We enforce these attribute names so later sessions can reuse your class!
        self.positive_label = "setosa"
        self.negative_label = "not_setosa"
\end{minted}

\textbf{Task:}
\begin{enumerate}
    \item Define this \texttt{\_\_init\_\_} method inside the \texttt{IrisRuleClassifier} class.
    \item In \texttt{main()}, create an instance: \texttt{my\_clf = IrisRuleClassifier(threshold=2.5)}.
    \item Print the stored threshold: \texttt{print("Threshold:", my\_clf.threshold)}.
\end{enumerate}
\end{tcolorbox}

\begin{tcolorbox}[title=Checkpoint]
You should be able to create two objects with different thresholds and see that each object stores its own value:
\begin{minted}[fontsize=\small, linenos, frame=lines]{python}
clf_a = IrisRuleClassifier(threshold=1.5)
clf_b = IrisRuleClassifier(threshold=2.5)
print(clf_a.threshold, clf_b.threshold)  # should print two different numbers
\end{minted}
\end{tcolorbox}
\begin{tcolorbox}[title=Common Mistake]
Forgetting \texttt{self}.
Inside methods you must write \texttt{self.threshold}, not just \texttt{threshold}, because \texttt{threshold} is not a global variable.
\end{tcolorbox}

\section{The Logic (\texttt{predict})}

\begin{tcolorbox}[title=Do: Implement predict]
In session 5, you wrote \texttt{classify\_sample(sample, settings)}.
In session 6, this becomes a \textbf{method}.
Because \texttt{threshold} and the labels already live on \texttt{self}, you no longer need a \texttt{settings} dictionary argument.

Conceptually, you are changing:
\texttt{classify\_sample(sample, settings)} $\rightarrow$ \texttt{my\_clf.predict(sample)}

\begin{minted}[fontsize=\small, linenos, frame=lines]{python}
    def predict(self, sample):
        # We assume sample is a dict with 'petal_length'
        if sample["petal_length"] < self.threshold:
            return self.positive_label
        return self.negative_label
\end{minted}

\textbf{Task:} Add this method to your class (indented inside the class). Identify which inputs are "per sample" and which inputs are "settings".
\texttt{sample} changes every iteration, so it stays an argument. \texttt{threshold} is a setting, so it becomes \texttt{self.threshold}.
\end{tcolorbox}

\begin{tcolorbox}[title=Checkpoint]
Call \texttt{my\_clf.predict(...)} on a single dictionary and make sure you get either \texttt{"setosa"} or \texttt{"not\_setosa"}.
\end{tcolorbox}
\begin{tcolorbox}[title=Common Mistake]
Forgetting to include \texttt{self} as the first parameter of a method. If you define \texttt{def predict(sample):}, Python will raise a TypeError when you call it!
\end{tcolorbox}

\section{Processing the Dataset (Evaluation)}

\begin{tcolorbox}[title=Do: Implement evaluate]
In session 5, you had a function \texttt{initialize\_predictions} that looped over the data.
We will now create a method \texttt{evaluate(self, dataset)} that does the entire prediction loop and calculates accuracy.

\textbf{Why make this a method?}
Because evaluation depends on \texttt{predict()}, and \texttt{predict()} depends on the object's configuration.
By keeping them together, you avoid passing settings around.

\textbf{Task:} Implement the \texttt{evaluate(self, dataset)} method. 
\begin{itemize}\setlength{\itemsep}{2pt}
    \item The evaluation loop should call \texttt{self.predict(sample)}.
    \item Determine truth (if species is "setosa", truth is \texttt{self.positive\_label}, else \texttt{self.negative\_label}).
    \item The method should keep counting logic local (\texttt{correct}, \texttt{total}).
    \item \textbf{CRITICAL:} Ensure \texttt{total += 1} runs for every sample, even if the prediction is wrong.
    \item Return accuracy (\texttt{correct / total} if \texttt{total > 0} else \texttt{0.0}).
\end{itemize}
\end{tcolorbox}

\begin{tcolorbox}[title=Checkpoint]
\begin{itemize}\setlength{\itemsep}{2pt}
    \item Accuracy is a number between 0.0 and 1.0.
    \item Test your method on a dataset with at least one misclassified sample and confirm that \texttt{total} equals the number of samples.
\end{itemize}
\end{tcolorbox}
\begin{tcolorbox}[title=Common Mistake]
Forgetting to increment \texttt{total} for every sample, even wrong ones. If you only increment total when correct, your accuracy will always be 100\% (which is wrong!).
\end{tcolorbox}

\section{Putting it all together (Main Script)}

\begin{tcolorbox}[title=Do: Use the Class in \texttt{main()}]
Now you can use your class to test different classifiers easily.
The workflow becomes:
\begin{enumerate}
    \item Use the provided \texttt{dataset} (list of dictionaries).
    \item Create one or more classifier objects (each with its own threshold).
    \item Ask each object to evaluate itself on the dataset.
\end{enumerate}

\begin{minted}[fontsize=\small, linenos, frame=lines]{python}
# Create Objects with different configurations
clf_strict = IrisRuleClassifier(threshold=1.5)   # Strict threshold
clf_loose  = IrisRuleClassifier(threshold=2.5)   # Loose threshold

# Evaluate each one
print("Strict Classifier Accuracy:", clf_strict.evaluate(dataset))
print("Loose Classifier Accuracy: ", clf_loose.evaluate(dataset))
\end{minted}

Notice that the dataset is passed in, but the threshold is not. The threshold lives inside each object!
\end{tcolorbox}

\section{Review and Next Steps}

\textbf{Summary:}
\begin{itemize}
    \item \textbf{Classes} group data (attributes) and behavior (methods) together.
    \item \textbf{Objects} remember their own configuration, allowing you to create many variations easily.
    \item Object-oriented design is the foundation of reusable software libraries. Machine-learning libraries like \texttt{scikit-learn} use objects to encapsulate models!
\end{itemize}

\textbf{Review Questions:}
\begin{enumerate}
    \item What does \texttt{self} refer to inside a method?
    \item Why did we move \texttt{threshold} from a dictionary into \texttt{self.threshold}?
    \item Why is inheritance useful? (Think ahead: what if we wanted a new classifier that guessed randomly, but shared the same evaluation logic?)
\end{enumerate}

\textbf{Micro-Exercise:} Create two classifier objects with different thresholds (e.g., 1.4 and 5.0) and compare their accuracies. Can you find a threshold that performs perfectly on the provided dataset?

\textbf{Stretch Exercise:} Write an additional method on the class called \texttt{predict\_all(self, dataset)} that returns a list of predictions for the entire dataset, rather than just returning accuracy.
